\chapter*{Abstract}

The Master's Thesis focuses on presenting the differences in usage of large language models, which belong to different configurations of RAG systems. The author has presented process and results of the research in terms of the algorithms' performance while answering questions related to external unstructured and structured data sources. To this end, the author has created a test workflow and an interface enabling conversation with the two chosen optimal systems.

The theory section introduces a reader to natural language processing, large language models and retrieval-augmented generation technique. There were described: history of chatbots, the way machines understand human language, how LLMs work and why they are eagerly used in RAG systems.

The chapter dedicated to literature review shows the scientists' approach to RAG systems evaluation. First reffered publication presented a scalable multi-document performance test called MEBench. In the second one, the authors delved into the topic of RGB benchmark, which enabled testing RAG systems with 1000 questions from four categories: noise robustness, negative rejection, information integration and counterfactual robustness. However, the third article focused on NL-to-SQL benchmark, which consisted of altogether 360 simple, moderate and challenging questions related to SQL database.

In the research part, the author has presented his original workflow enabling test automation of thirty different RAG systems. The creation process of unstructured data (PDF and text file) and structured data (SQL database and CSV file) was also described. Inspired by studied literature, the author proposed two sets of eight questions dedicated to two of the mentioned types of information sources.

The last section focuses on results analysis and presents comparison of RAG systems in terms of accuracy and cost of generated answers, time of generation and token throughput. In the last step, the two most efficient systems were chosen, one from each category. Then, the author has integrated them with a graphical user interface, enabling interaction through conversation.



\bigskip
\bigskip

\noindent \textbf{Key words}: \\
NLP, Chatbot, artificial intelligence, LLM, RAG, AI agent, unstructured data, structured data, PDF, text file, SQL, CSV, automation, GUI, scraping, embedder, vector database


\bigskip

\noindent \textbf{Field of science and technology as required by the OECD}: \\
Engineering and Technical Sciences, Technical Computer Science and Telecommunications